\section{Limitations and conclusion}\label{sec:conclusion}

\paragraph{Limitations.} Seven are material. First, the residual: our
method removes most of the conditional coverage deficit, not all of it:
88.1\% against 90\% nominal in stress, with the shortfall concentrated in
the onset window of Section~\ref{sec:onset}, and 99\% VaR stress breaches
of 1.8\% against 1\%. We consider reporting this residual preferable to
tuning until a table reads 90.0 on a slice chosen after the fact. Second,
scope: one market (US large-cap equities), one target (daily realized
volatility); the mechanism (rare-regime starvation repaired by
cross-sectional pooling) should transfer to any panel of related
series, but we have not shown it. Third, the guarantees control realized
frequencies in regime-weighted time, conditional on the estimated
memberships; they are not day-level conditional probability statements,
which are unattainable without distributional assumptions
\citep{foygelbarber2021limits}, and the hindsight ablation bounds, but
does not eliminate, the gap to latent-truth conditioning. Fourth,
although the point-in-time universe removes survivorship in stock
selection, the realized-volatility source begins in 2003 and large-cap
liquid names are the easiest case for high-frequency measurement; broader
cross-sections deserve a dedicated study. Fifth, the finite-sample
constants: evaluated at the observed stress sample the per-regime bound
of Proposition~\ref{prop:adaptive} is vacuous
(Remark~\ref{rem:numeric}); the theorems supply worst-case pathwise
control, and the empirical tables, not the constants, carry the
finite-sample evidence. Sixth, the temporal analysis slices causally
issued intervals by subperiod; because the method has no fitting step
this is the frozen-design evaluation, but it is not a holdout kept
untouched during development. Seventh, the \emph{average} stress gain of
regime conditioning over pooled $K = 1$ is small and not statistically
significant; the case for regimes rests on the conditional structure
--- transition coverage, VaR balance, per-regime guarantees --- which we
quantify with dedicated uncertainty estimates rather than assert.

\paragraph{Conclusion.} Uncertainty statements attached to volatility
forecasts are consumed conditionally, on the bad days, and this
paper shows that the standard way of producing them can be valid only on
average: over a 2010--2025 evaluation of two hundred companies, online
conformal intervals promised at 90\% deliver 83\% on stress-regime days.
The deficit is not uniform history: it is severe in the fast regime
shifts of 2018--2025 (78.8\% stress coverage) and largely absent in the
slower-building episodes of 2010--2017 (89.6\%), and foundation-model
uncertainty fails in the same direction. Marginal validity conceals a
conditional failure mode of the recent market era, which is exactly when
such intervals are consumed. A calibration layer that pools the cross-section within
online-estimated regimes and tunes itself removes most of that deficit at
no detected cost in average interval score, carries finite-horizon
pathwise bounds on aggregate realized per-regime miss frequencies,
robust to arbitrary dependence, and maps directly to VaR,
where it is the only method compared whose calm and stress breach rates
are statistically balanced at the 5\% level, with top-tier backtest
pass shares. What
it does not do by default (cover the second day of a volatility spike)
recurs across every method, era, and conditioning scheme we evaluate,
though the estimate rests on few observed transitions and covering the
window is a priced trade rather than an impossibility
(Section~\ref{sec:onset-pareto}). The practical reading for forecasters is
simple: report coverage conditionally, expect marginal guarantees to hide
exactly the failures that matter, and pool calibration information across
the cross-section when rare states are the concern.

\paragraph{Reproducibility.} All results regenerate from a single driver
script with fixed seeds; data sources are public or standard academic
subscriptions (CRSP via WRDS; Risk Lab realized volatilities; FRED/CBOE
state variables), with every data query logged. Code and full replication
instructions accompany the submission.
